\documentclass[10pt,a4paper]{article}
\usepackage[utf8]{inputenc}
\usepackage[T1]{fontenc}
\usepackage[spanish,es-nodecimaldot]{babel}
\usepackage{amsmath,amssymb}
\usepackage{newtxtext,newtxmath}
\usepackage{microtype}
\usepackage[margin=1.8cm,top=2.2cm,bottom=2.2cm]{geometry}
\usepackage{booktabs,tabularx,colortbl}
\usepackage{xcolor}
\usepackage{fancyhdr}
\usepackage{enumitem}
\usepackage{titlesec}
\usepackage[skins,breakable]{tcolorbox}
\usepackage[colorlinks=true,linkcolor=embatnavy,citecolor=embatnavy,urlcolor=accentblue]{hyperref}

% ─── Paleta de colores corporativa y técnica ────────────────────────────────
\definecolor{embatnavy}{HTML}{050B2C}
\definecolor{embatteal}{HTML}{00A896}
\definecolor{accentblue}{HTML}{1A56DB}
\definecolor{boxbg}{HTML}{F8FAFC}
\definecolor{boxborder}{HTML}{CBD5E1}
\definecolor{textmuted}{HTML}{475569}
\definecolor{alertred}{HTML}{991B1B}

% ─── Configuración de cabeceras y pie de página ──────────────────────────────
\pagestyle{fancy}
\fancyhf{}
\fancyhead[L]{\small\color{textmuted}\textbf{HackSpain 2026} $\cdot$ Reto Embat (X-Ray) $\cdot$ Technical Whitepaper}
\fancyhead[R]{\small\color{textmuted}Septiembre 2026}
\fancyfoot[C]{\small\color{textmuted}Página \thepage}
\renewcommand{\headrulewidth}{0.4pt}
\renewcommand{\footrulewidth}{0.4pt}

% ─── Formato de títulos de sección ──────────────────────────────────────────
\titleformat{\section}{\large\bfseries\color{embatnavy}}{\thesection.}{0.5em}{}[\titlerule]
\titleformat{\subsection}{\normalsize\bfseries\color{embatnavy}}{\thesubsection.}{0.5em}{}
\titlespacing*{\section}{0pt}{1.2em}{0.6em}
\titlespacing*{\subsection}{0pt}{0.8em}{0.4em}

% ─── Cajas de estilo ─────────────────────────────────────────────────────────
\tcbset{
  paperbox/.style={
    enhanced,
    colback=boxbg,
    colframe=boxborder,
    arc=2mm,
    boxrule=0.7pt,
    left=3mm, right=3mm, top=2.5mm, bottom=2.5mm,
    fonttitle=\bfseries\color{embatnavy}
  },
  formulabox/.style={
    enhanced,
    colback=embatnavy!3!white,
    colframe=embatnavy,
    arc=2mm,
    boxrule=1pt,
    left=4mm, right=4mm, top=3mm, bottom=3mm
  }
}

\begin{document}

% ─── Encabezado Principal del Paper ──────────────────────────────────────────
\begin{center}
  {\fontsize{18}{22}\selectfont\bfseries\color{embatnavy} X-Ray: Motor de Scoring Continuo, Evaluación de Trayectorias Financieras\\ y Monitor Temprano de Estrés de Tesorería B2B}\\[0.4em]
  {\large\color{textmuted} Especificación matemática, desacoplamiento de calidad documental y validación empírica en 1.286 sociedades}\\[0.8em]
  
  {\small\textbf{Antonio Machuca} \quad $\cdot$ \quad \textbf{Carlos León} \quad $\cdot$ \quad \textbf{Hugo Oliva} \quad $\cdot$ \quad \textbf{Pedro Jon} \quad $\cdot$ \quad \textbf{Agustín Muñoz}}\\[0.2em]
  {\footnotesize\color{textmuted} HackSpain 2026 $\cdot$ Track Embat (X-Ray) $\cdot$ Escuela Técnica Superior de Ingenieros de Telecomunicación (UPM-ETSIT), Madrid}\\[0.2em]
  {\footnotesize\color{accentblue}\texttt{https://github.com/antoniomachuca/hackspain-2026-embat}}
\end{center}

\vspace{0.3em}

% ─── Resumen Ejecutivo / Abstract ────────────────────────────────────────────
\begin{tcolorbox}[paperbox, title=Resumen Ejecutivo y Declaración de Alcance Metodológico]
\small
\textbf{Problema analítico:} La evaluación de solvencia empresarial basada exclusivamente en fotos contables fijas u agregados estáticos oculta trayectorias dinámicas contrapuestas. Una compañía con un nivel base moderado en trayectoria ascendente puede presentar menor riesgo crediticio y operativo que una sociedad históricamente sólida pero en fase de deterioro acelerado. Asimismo, la ausencia o baja densidad de registros en entornos pyme suele contaminar las estimaciones financieras, confundiéndose la falta de datos con solvencia deficiente o riesgo nulo.

\textbf{Metodología formulada:} Presentamos la especificación matemática y validación econométrica de \textbf{X-Ray}, un motor continuo y determinista acotado estrictamente en $[0, 100]$. El modelo introduce cuatro principios de diseño formal:
\begin{enumerate}[leftmargin=*, noitemsep, topsep=1pt]
  \item \textbf{Desacople estricto entre solvencia y gobernanza documental:} Se suprime la contracción ciega de pilares hacia valores neutrales ($0{,}5$). La calidad de datos se aísla en un indicador ortogonal independiente, el \emph{Data Confidence Index} ($\text{DCI}_t = 100 \cdot \overline{q}_t^{[3]} \cdot \min(m/6, 1{,}0)$), reduciendo la correlación espuria con calidad documental a $\rho = 0{,}107$ en el corte final (máximo histórico $\rho = 0{,}211 \ll 0{,}65$).
  \item \textbf{Crecimiento sin estrangulación:} Sustitución del promedio geométrico por el promedio aritmético en el componente de crecimiento bancario ($G_t^b$), preservando el reconocimiento de fases tempranas de recuperación operativa.
  \item \textbf{Estrés comercial y fragilidad desprovista de atenuación artificial:} Eliminación de factores multiplicativos de cobertura en la penalización por fragilidad ($F_t$) e integración del estrés de facturas vencidas ($\text{invoice\_stress}$).
  \item \textbf{Monitor proactivo con corte de histéresis (*opposing momentum breaker*):} Filtro de persistencia trimestral no solapado e interrupción automática de estados positivos ante shocks adversos ($M_t < -0{,}10$ o $\Delta S_t < -3{,}0$), erradicando alertas invertidas en series reales.
\end{enumerate}

\textbf{Evidencia experimental cuantificada:} En cohortes sintéticas estandarizadas (400 series por escenario), el recall a $\le 6$ meses alcanza el 94,0\% en deterioro gradual frente al 92,5\% en mejora y 91,5\% en recuperación (gap simétrico de 1,5--2,5 p.p.; retraso mediano idéntico de 4,0 meses). En el benchmark de agotamiento de caja, el 100\% de las series en deterioro gradual son detectadas antes de la insolvencia (antelación mediana de 14,0 meses, P10 = 8,0 meses), con el 99,0\% de las alertas emitidas con un nivel base aún alto ($B_t \ge 60$). Ante shocks puntuales (\emph{pulse}), la tasa de falsas alarmas es estrictamente 0,000 por empresa-año. La evaluación en la partición reservada del dataset real (50 grupos, 239 sociedades, 365 cortes evaluables) arroja independencia técnica exacta de lote ($\Delta_{\text{batch}} = 0{,}000$).
\end{tcolorbox}

\vspace{0.2em}

% ─── Sección 1: Correspondencia con las 6 preguntas de Embat ─────────────────
\section*{1. Correspondencia Sistemática con las Seis Preguntas del Reto}

El marco analítico de X-Ray proporciona una respuesta matemática y operativa a cada uno de los seis requerimientos planteados en el reto Embat:

\begin{center}
\small
\begin{tabularx}{\textwidth}{p{2.6cm}p{3.2cm}p{4.2cm}X}
\toprule
\textbf{Requerimiento} & \textbf{Concepto Econométrico} & \textbf{Mecanismo en el Motor} & \textbf{Validación Empírica y Límites} \\
\midrule
\textbf{1. Quién está sano} & Solvencia de flujo y cobertura corriente & Nivel base ponderado $B_t$ y bandas de salud (\texttt{SOLIDA}, \texttt{INTERMEDIA}, \texttt{DEBIL}) & Score acotado en $[0, 100]$. Desacoplado de la cobertura documental. No presupone inferencia de default externo sin etiquetas de supervisión. \\
\textbf{2. Quién mejora} & Expansión sostenida y trayectoria positiva & Momentum $M_t > 0{,}10$, crecimiento $G_t$ y estados \texttt{MEJORANDO} y \texttt{RECUPERACION} & Recall sintético $\le 6$ meses del 92,5\% en mejora y 91,5\% en recuperación. Retraso mediano de 4,0 meses sin retrotraer fechas. \\
\textbf{3. Quién se tuerce} & Deterioro temprano con base aún sólida & Momentum negativo confirmado ($M_t < -0{,}10$) sobre base $B_t \ge 60$ (\texttt{TORCIENDOSE}) & 99,0\% de las alertas de deterioro gradual se emiten con base $B_t \ge 60$. Alerta antes de que el daño patrimonial sea irreversible. \\
\textbf{4. Bache vs. caída} & Shock estocástico transitorio frente a tendencia & Medianas móviles trimestrales, 3 meses de persistencia y estado \texttt{BACHE} & Tasa de falsas alertas de 0,000 por empresa-año en shocks de un mes (\emph{pulse}). Mitigación del 96\% en ciclos estacionales con contexto interanual. \\
\textbf{5. Por qué cambia} & Atribución aditiva exacta de variaciones & Descomposición lineal término a término (Waterfall) con corrección de clipping & Identidad aritmética exacta: $\Delta S_t = \sum_k \Delta \text{Points}_k + \Delta_t^{\mathrm{clip}}$. Residuo algebraico nulo ($\|\epsilon\| < 10^{-13}$). \\
\textbf{6. Cuándo se avisó} & Antelación temporal a la insolvencia de caja & Demonio proactivo \texttt{score\_monitor.py} con cursor temporal inmutable & Antelación mediana de 14,0 meses al agotamiento de caja en deterioro gradual (7 a 8 meses en control determinista). Lead nulo en feed si no hay evento fechado. \\
\bottomrule
\end{tabularx}
\end{center}

% ─── Sección 2: Arquitectura Matemática y Ecuaciones Actualizadas ───────────
\section*{2. Arquitectura Matemática y Ecuación Global del Score}

El índice sintético de salud financiera $S_t$ se define para cada corte temporal $t$ mediante una combinación linealmente separable acotada mediante el operador de recorte (\emph{clipping}):

\begin{tcolorbox}[formulabox]
\vspace{-0.2em}
\[
\boxed{S_t = \operatorname{clip}_{[0,100]}\left( B_t + \lambda M_t + \gamma G_t - \delta F_t \right)}
\]
\[
B_t = 100 \left( w_L L_t + w_C C_t + w_D D_t \right), \qquad \sum_{k \in \{L, C, D\}} w_k = 1, \quad w_k \geq 0
\]
\vspace{-0.4em}
\end{tcolorbox}

Donde $L_t, C_t, D_t, G_t, F_t \in [0, 1]$, el nivel base $B_t \in [0, 100]$, y el vector de momentum $M_t \in [-1, 1]$. La parametrización obedece a criterios de prudencia financiera y estabilidad econométrica:

\begin{center}
\small
\begin{tabularx}{\textwidth}{lcX}
\toprule
\textbf{Parámetro} & \textbf{Valor} & \textbf{Fundamentación Económica y Decisión de Diseño} \\
\midrule
$w_L$ & 0,50 & Ponderación de liquidez y cobertura operativa corriente: capacidad inmediata de absorber compromisos. \\
$w_C$ & 0,30 & Eficiencia de cobros bancarios (baja tasa de devoluciones y regularidad), enriquecible opcionalmente con ERP. \\
$w_D$ & 0,20 & Carga de servicio de deuda bancaria observada. Se mantiene conservador ante falta de calendarios integrales. \\
$\lambda$ & 8 & Amplitud máxima de ajuste por trayectoria en $[-8, +8]$ puntos. Permite a empresas en mejora superar a bases estáticas superiores. \\
$\gamma$ & 6 & Bonificación máxima por crecimiento de calidad validado mediante conversión en flujo de caja positivo. \\
$\delta$ & 8 & Penalización máxima por fragilidad operativa, desajustes de timing intrames y concentración de contrapartes. \\
\bottomrule
\end{tabularx}
\end{center}

A lo largo del chasis matemático se emplea recurrentemente la función saturante de tipo Hill / Michaelis-Menten:
\[
h(x; a) = \frac{\max(x, 0)}{\max(x, 0) + a}, \quad a > 0
\]
Dicha función garantiza monotonicidad estricta, diferenciabilidad continua en $(0, \infty)$, acotación exacta en $[0, 1)$, derivada acotada por $1/a$, y un punto de calibración equidistante $h(a; a) = 0{,}5$.

% ─── Sección 3: Identificabilidad y Desacople de Calidad Documental ──────────
\section*{3. Identificabilidad Causal y Desacoplamiento de Calidad Documental}

\subsection*{3.1. Identificabilidad point-in-time y no retropropagación}
El conjunto de datos procesado comprende 1.286 sociedades agrupadas en 250 grupos empresariales con cortes mensuales desde octubre de 2024 hasta septiembre de 2026. La modelización causal respeta estrictamente dos restricciones de identificabilidad:
\begin{enumerate}[leftmargin=*, noitemsep]
  \item \textbf{Ausencia de ancla de caja histórica:} El saldo de caja al instante $t$ satisface $K_t = K_0 + \sum_{u \leq t} f_u$, donde $f_u$ denota el flujo neto de movimientos bancarios. Sin conocer el saldo inicial $K_0$, existen infinitas trayectorias de tesorería compatibles con los flujos observados. El archivo \texttt{balances.csv} suministra únicamente la posición a 01/09/2026. Reconstruir retrospectivamente $K_t$ hacia el pasado a partir de dicha foto constituiría una violación causal grave (\emph{lookahead leakage}). Por ello, el núcleo histórico del motor evalúa la cobertura operativa de flujos sin asumir saldos pasados.
  \item \textbf{Cero observaciones conocidas $\neq$ saldo nulo:} En el núcleo bancario universal, el número de observaciones de caja histórica verificada es exactamente cero. Imputar un saldo nulo distorsionaría la medición; por ende, $L_t$ se calcula mediante ratios de flujo.
\end{enumerate}

\subsection*{3.2. Desacoplamiento formal de la solvencia frente a la cobertura de datos}
En formulaciones iniciales, los factores de solvencia se contraían hacia un prior neutral ($0{,}5$) mediante un operador ciego de atenuación: $\mathcal{A}(x, \rho_t) = 0{,}5 + \rho_t (x - 0{,}5)$. Este diseño presentaba un defecto econométrico grave: introducía un sesgo de selección artificial donde la falta de densidad transaccional castigaba a sociedades solventes y, simétricamente, amortiguaba el riesgo de sociedades en insolvencia incipiente.

En la arquitectura actual, \textbf{se elimina completamente la contracción ciega en los tres pilares del nivel base} ($L_t, C_t, D_t$). Cada pilar refleja fielmente la calidad intrínseca de los flujos observados. Para garantizar la gobernanza de datos y la transparencia analítica para el suscriptor de riesgos, la disponibilidad y madurez de la información se aíslan en una variable ortogonal:

\begin{tcolorbox}[formulabox]
\vspace{-0.2em}
\[
\boxed{\text{DCI}_t = 100 \cdot \overline{q}_t^{[3]} \cdot \min\left(\frac{m_t}{6}, 1{,}0\right)}
\]
\vspace{-0.4em}
\end{tcolorbox}

Donde $\overline{q}_t^{[3]} \in [0, 1]$ es la cobertura bancaria media del último trimestre y $m_t$ representa el número de meses con transacciones observadas en el historial. El \emph{Data Confidence Index} ($\text{DCI}_t \in [0, 100]$) se reporta de forma independiente junto al score, informando la fiabilidad de la muestra sin alterar el diagnóstico de solvencia. 

Empíricamente, este desacoplamiento erradica la dependencia espuria: la correlación de Spearman entre el score final y la calidad de datos en el mes 23 es de tan solo $\rho = 0{,}107$, con un máximo en todos los cortes maduros de $\rho = 0{,}211$, situándose muy por debajo del umbral regulatorio de colinealidad de $0{,}65$.

% ─── Sección 4: Pilares del Nivel Base: Cobertura, Cobros y Deuda ────────────
\section*{4. Pilares del Nivel Base: Cobertura Operativa, Cobros y Endeudamiento}

Para cada mes $t$, se definen los cobros netos $R_t \ge 0$, los pagos operativos $E_t \ge 0$, el servicio de deuda registrado $H_t \ge 0$, y sus sumas móviles trimestrales $R_t^{[3]}, E_t^{[3]}, H_t^{[3]}$.

\subsection*{4.1. Cobertura operativa y liquidez ($L_t$)}
En ausencia de saldo de caja auditado, el pilar $L_t^b$ mide el margen relativo entre cobros y pagos corrientes en ventana de 3 meses:
\[
u_t = \frac{R_t^{[3]} - E_t^{[3]}}{R_t^{[3]} + E_t^{[3]}}, \qquad L_t^b = \begin{cases} 0{,}5 + 0{,}5 \tanh\left(\frac{u_t}{0{,}50}\right), & \text{si } R_t^{[3]} + E_t^{[3]} > 0, \\ 0{,}5, & \text{sin evidencia.} \end{cases}
\]
La función asigna exactamente $0{,}5$ en equilibrio operativo ($R = E$). Si en un corte específico se aporta saldo de caja certificado $K_t \in \mathbb{R}$, con $K_t^+ = \max(K_t, 0)$, tasa mensual de consumo de caja (\emph{burn rate}) $b_t = \max\left(\frac{E_t^{[3]} + H_t^{[3]} - R_t^{[3]}}{3}, 0\right)$ y pasivo corriente a corto plazo $P_t^*$:
\[
\ell_t^{\mathrm{run}} = \frac{K_t^+}{K_t^+ + 3b_t}, \qquad \ell_t^{\mathrm{cov}} = \frac{K_t^+}{K_t^+ + P_t^*}, \qquad L_t = 0{,}5 \ell_t^{\mathrm{run}} + 0{,}5 \ell_t^{\mathrm{cov}}.
\]

\subsection*{4.2. Calidad de cobro y convergencia ERP ($C_t$)}
El pilar bancario pondera a partes iguales la penalización por devoluciones y la regularidad del flujo de entradas:
\[
f_t = \frac{U_t^{[3]}}{(R^{\mathrm{bruto}})_t^{[3]}}, \quad s_t^{\mathrm{ref}} = 1 - h(f_t; 0{,}05), \quad s_t^{\mathrm{reg}} = \frac{1}{1 + CV_t^{[6]}}, \quad C_t^b = 0{,}5 s_t^{\mathrm{ref}} + 0{,}5 s_t^{\mathrm{reg}},
\]
donde $CV_t^{[6]}$ es el coeficiente de variación de los cobros en 6 meses. Cuando se dispone de integración ERP fechada, con $DSO_t$, variación trimestral $\Delta DSO_t = DSO_t - DSO_{t-3}$ y proporción de facturas vencidas impagadas $p_t^{\mathrm{late}}$:
\[
z_t = 1 - h(DSO_t; 60), \quad z_t^* = 0{,}75 z_t + 0{,}25 \left(0{,}5 - 0{,}5 \tanh\left(\frac{\Delta DSO_t}{30}\right)\right), \quad C_t^e = 0{,}5 z_t^* + 0{,}5(1 - p_t^{\mathrm{late}}).
\]
El pilar enriquecido adopta la mezcla convexa $C_t = (1 - \eta_t) C_t^b + \eta_t C_t^e$, con $\eta_t = 0{,}40 q_t^e \in [0, 0{,}40]$.

\subsection*{4.3. Carga de endeudamiento financiero ($D_t$)}
Evaluamos la fracción de cobros absorbida por amortizaciones de principal e intereses:
\[
d_t^H = \frac{H_t^{[3]}}{R_t^{[3]} + H_t^{[3]}}, \qquad D_t = 0{,}5 - 0{,}5 \tanh\left(\frac{d_t^H}{0{,}25}\right).
\]
Por diseño conservador, $D_t \in [0, 0{,}5]$: una mayor carga de deuda reduce monótonamente la puntuación, pero la ausencia de pagos de deuda observados conserva el prior neutral $0{,}5$, impidiendo calificar a la sociedad como desapalancada sin inventario crediticio verificado.

% ─── Sección 5: Crecimiento de Calidad y Fragilidad Comercial ─────────────────
\section*{5. Crecimiento de Calidad y Fragilidad Comercial Integrada}

\subsection*{5.1. Crecimiento de calidad ($G_t$) y justificación del promedio aritmético}
El componente de crecimiento premia la expansión de cobros que se traduce efectivamente en tesorería, modulado por la cobertura temporal:
\[
g_t^b = \frac{R_t^{[3]} - R_{t-3}^{[3]}}{R_t^{[3]} + R_{t-3}^{[3]}}, \qquad G_t^b = \min(\rho_t, \rho_{t-3}) \cdot \tanh\left(\frac{\max(g_t^b, 0)}{0{,}20}\right) \cdot \left[0{,}5 L_t^b + 0{,}5 C_t^b\right].
\]
\textbf{Fundamentación analítica frente a la media geométrica:} La especificación previa empleaba el término multiplicativo $\sqrt{L_t^b C_t^b}$. En situaciones de recuperación operativa donde una empresa comienza a salir de un shock adverso, los ratios de liquidez pasada pueden situarse en niveles moderadamente bajos mientras que la regularidad de cobro ya se ha restablecido. La media geométrica sufre un fenómeno de \emph{estrangulación analítica} ($\lim_{x \to 0} \sqrt{x \cdot y} = 0$), anulando por completo el incentivo de crecimiento y retrasando la detección del giro positivo. El promedio aritmético $[0{,}5 L_t^b + 0{,}5 C_t^b]$ actúa como un amortiguador equilibrado que preserva la señal de recuperación sin premiar crecimiento tóxico (donde ambos factores colapsarían).

\subsection*{5.2. Fragilidad operativa y estrés de facturación ($F_t$)}
La fragilidad evalúa desajustes intrames de timing, saldos negativos e impagos comerciales:
\[
F_t = (1 - \theta_t) s_t + \theta_t h(\text{HHI}_t; 0{,}25), \qquad \theta_t = 0{,}50 \min(q_t^{\mathrm{HHI}}, 1).
\]
\textbf{Corrección del factor de cobertura:} Se suprime la multiplicación previa por la cobertura bancaria ($\rho_t F_t$). Dicha formulación reducía artificialmente la fragilidad estimada cuando empeoraba la calidad del dato, incentivando la opacidad informativa. 

El estrés combinado $s_t$ se formula mediante complementos de supervivencia estocástica:
\[
s_t = 1 - (1 - s_t^{\mathrm{timing}})(1 - s_t^{\mathrm{neg}})(1 - s_t^{\mathrm{late}}),
\]
donde $s_t^{\mathrm{timing}} = h\left(\frac{\overline{\mathrm{gap}}_t^{[3]}}{(E_t^{[3]}+H_t^{[3]})/3}; 0{,}50\right)$, $s_t^{\mathrm{neg}} = h(\text{fracción intrames con saldo } < 0; 0{,}25)$, y el estrés comercial por facturas vencidas:
\[
s_t^{\mathrm{late}} = h(p_t^{\mathrm{late}}; 0{,}25) \cdot \mathbf{1}_{\{q_t^e > 0\}}.
\]
La integración de $s_t^{\mathrm{late}}$ permite detectar estrés financiero comercial semanas o meses antes de que se refleje en descubiertos bancarios.

% ─── Sección 6: Análisis Formal y Empírico: "Si Acierta" ────────────────────
\section*{6. Análisis Formal y Empírico: Bloque <<Si Acierta>>}

\subsection*{6.1. Generalización e independencia técnica de lote}
Para verificar que el motor no introduce fugas de información ni acoplamientos espurios entre empresas procesadas en el mismo lote, se ejecutó la prueba de permutación y partición de lote:
\[
\Delta_{\text{batch}} = \max_{i, t} |S_{i, t}^{\text{aislado}} - S_{i, t}^{\text{lote}}| = 0{,}000.
\]
La diferencia idénticamente nula certifica que el cómputo para cada sociedad es estrictamente invariante respecto al conjunto de empresas acompañantes.

\subsection*{6.2. Evaluación en holdout real y transparencia sobre la cobertura futura}
Se fijó una partición reservada (\emph{holdout}) de 50 grupos empresariales (semilla 42), abarcando 239 sociedades. Tras exigir los criterios de madurez del monitor (\texttt{minimum\_quality = 0.50}, ventanas trimestrales completas), restan 365 observaciones empresa-corte utilizables distribuidas en 199 sociedades de 41 grupos.

\begin{center}
\small
\begin{tabularx}{\textwidth}{lXcc}
\toprule
\textbf{Predictor evaluado} & \textbf{Concepto analítico} & \textbf{Spearman ($r_s$)} & \textbf{Intervalo bootstrap 95\%} \\
\midrule
Score dinámico $S_t$ & Solvencia integral y trayectoria de flujo & -0,062 & $[-0{,}206; +0{,}087]$ \\
Cobertura histórica $\rho_t^{[3]}$ & Calidad y densidad de datos a 3 meses & -0,087 & $[-0{,}277; +0{,}107]$ \\
Cobertura histórica $\rho_t^{[6]}$ & Calidad y densidad de datos a 6 meses & -0,037 & $[-0{,}235; +0{,}140]$ \\
\bottomrule
\end{tabularx}
\end{center}

Todos los intervalos bootstrap agrupados por conglomerados (41 grupos, 300 réplicas) cruzan holgadamente el cero. Esto confirma un principio de honestidad econométrica: el score mide solvencia contable y regularidad de flujos, pero no predice el volumen futuro de transacciones bancarias mejor que una persistencia histórica de seis meses ($q_6$). El sistema reporta estos hallazgos sin ajustar artificialmente los hiperparámetros para sobreajustar proxies auxiliares.

\subsection*{6.3. Dinámica de trayectoria y el mecanismo anti-enganche (*opposing momentum breaker*)}
El vector de momentum $M_t$ sintetiza la derivada de la base bancaria $v_t = (B_t^b - B_{t-3}^b)/3$ y el diferencial de medias móviles exponenciales rápida y lenta ($E_{3,t} - E_{6,t}$), filtrado por el coeficiente de persistencia direccional no solapado $c_t \in [0, 1]$:
\[
M_t = c_t \left[ 0{,}5 \tanh(v_t / 2{,}0) + 0{,}5 \tanh((E_{3,t} - E_{6,t}) / 5{,}0) \right].
\]
Para clasificar los estados financieros (\texttt{ESTABLE}, \texttt{TORCIENDOSE}, \texttt{DETERIORO}, \texttt{MEJORANDO}, \texttt{RECUPERACION}), el clasificador exige 3 meses consecutivos de momentum fuera de $\pm 0{,}10$. Con el fin de evitar fluctuaciones espurias, una regla de histéresis retiene el estado previo durante 2 meses neutros (\texttt{hold}).

\textbf{El problema de enganche de estado:} En versiones preliminares, si una empresa en estado positivo (\texttt{MEJORANDO} o \texttt{RECUPERACION}) sufría un shock negativo repentino, la regla de histéresis mantenía el estado positivo durante los dos meses neutros de transición, emitiendo alertas invertidas absurdas durante fases de desplome.

\textbf{Solución implementada (*opposing momentum breaker*):} Se incorporan guardas deterministas de ruptura inmediata:
\[
\text{opposing\_negative} = \mathbf{1}_{\{S_{t-1} \in \text{POSITIVE\_STATES}\}} \land \left( M_t < -0{,}10 \;\lor\; \Delta S_t < -3{,}0 \right),
\]
\[
\text{opposing\_positive} = \mathbf{1}_{\{S_{t-1} \in \text{NEGATIVE\_STATES}\}} \land \left( M_t > +0{,}10 \;\lor\; \Delta S_t > +3{,}0 \right).
\]
Cuando se detecta un shock en sentido contrario, se anula la retención por histéresis y la sociedad transita inmediatamente a \texttt{ESTABLE}.

\textbf{Caso empírico de validación: \texttt{COMP\_0010}:} En el corte \texttt{2026-02-01}, la sociedad se hallaba en \texttt{RECUPERACION} con score $S = 69{,}86$, base $B = 64{,}88$ y momentum $M = +0{,}775$. En el corte posterior (\texttt{2026-03-01}), sufrió un desplome financiero abrupto, cayendo el score a $S = 39{,}00$ ($\Delta S = -30{,}86$ puntos) y el momentum a $M = -0{,}471$. Con el mecanismo anti-enganche, el estado positivo se desactivó en el instante exacto del impacto, pasando a \texttt{ESTABLE} e impidiendo la emisión de alertas de recuperación durante la caída.

\subsection*{6.4. Las dos caras: simetría bidireccional en mejora y deterioro}
Un motor de scoring riguroso debe responder con igual sensibilidad y presteza ante recuperaciones que ante deterioros. En la evaluación de cohortes sintéticas estandarizadas (400 series por tipología):
\begin{itemize}[leftmargin=*, noitemsep]
  \item \textbf{Recall a $\le 6$ meses:} 94,0\% en deterioro gradual vs. 92,5\% en mejora sostenida y 91,5\% en recuperación operativa (gap simétrico de tan solo 1,5 a 2,5 p.p.).
  \item \textbf{Celeridad temporal:} El retraso mediano de detección es estrictamente idéntico: \textbf{4,0 meses} en las tres dinámicas.
\end{itemize}
Esto demuestra empíricamente que la introducción de la media aritmética en $G_t^b$ y la simetría del estimador de momentum logran una respuesta bidireccional equilibrada sin sesgos conservadores asimétricos.

% ─── Sección 7: Análisis Formal y Empírico: "Si Llega a Tiempo" ─────────────
\section*{7. Análisis Formal y Empírico: Bloque <<Si Llega a Tiempo>>}

\subsection*{7.1. Anticipación cuantitativa del colapso de tesorería}
Para medir con rigor cuántos meses antes avisa el motor antes de que el daño sea evidente, se evaluaron 400 trayectorias de estrés de caja donde los flujos entran en déficit continuo a partir del mes 12 y se audita el mes en que el saldo acumulado se vuelve no positivo:
\begin{enumerate}[leftmargin=*, noitemsep]
  \item \textbf{Detección previa al agotamiento:} En deterioro gradual, el \textbf{100\%} de las empresas que colapsan (353 de 353) son alertadas antes de la quiebra de caja. La antelación mediana es de \textbf{14,0 meses} (con un percentil 10 de \textbf{8,0 meses} de margen mínimo garantizado).
  \item \textbf{Deterioro abrupto:} En shocks súbitos de alta severidad, el \textbf{88,2\%} (352 de 399 colapsos) se detectan antes de la pérdida total de tesorería, con una antelación mediana de \textbf{6,0 meses}.
  \item \textbf{Alerta con base todavía sana:} El \textbf{99,0\%} de las detecciones en deterioro gradual se producen cuando el nivel base $B_t$ es todavía superior o igual a 60, emitiendo el estado temprano \texttt{TORCIENDOSE} antes de que la foto patrimonial refleje insolvencia.
  \item \textbf{Control determinista de tesorería:} En el escenario de control estándar (caja inicial 300, cobros 120 que caen a 5 tras el mes 12), el monitor operativo emite su primera alerta confirmada en el mes 16, mientras que la caja se agota en el mes 23: \textbf{7 meses de antelación neta} sin utilizar la caja como variable explicativa.
\end{enumerate}
\textbf{Delimitación estricta:} Esta antelación está validada sobre el agotamiento del saldo de tesorería y facturas impagadas en el simulador causal, no sobre declaraciones mercantiles de concurso que no constan en los CSV del concurso.

\subsection*{7.2. Estabilidad frente a ruido y discriminación de shocks transitorios}
El algoritmo debe distinguir con precisión un bache puntual de una crisis estructural de liquidez:
\begin{itemize}[leftmargin=*, noitemsep]
  \item \textbf{Escenario de shock puntual (\emph{pulse}):} En 400 series sometidas a una caída brusca de márgenes de flujo de un solo mes, la tasa de falsas alertas adversas es exactamente \textbf{0,000 por empresa-año} (0,0\% de series afectadas). El requisito de 3 meses de persistencia temporal y la asignación del estado transitorio provisional \texttt{BACHE} absorben completamente el impacto sin disparar avisos ejecutivos.
  \item \textbf{Escenario estable con ruido estocástico:} Tasa de falsas alarmas de \textbf{0,000 por empresa-año}.
  \item \textbf{Escenario estacional cíclico:} La incorporación del contexto interanual ($t \ge 14$, comparando ventanas trimestrales con las homólogas del año precedente bajo una banda muerta de 0,06) reduce la tasa de falsas alarmas anuales de 0,215 a tan solo \textbf{0,0123 alertas por empresa-año}, lo que representa una mitigación del \textbf{96\%} del ruido estacional.
\end{itemize}

\subsection*{7.3. Monitor proactivo autónomo y cobertura de cartera}
El monitor no espera a que un analista consulte la posición de una compañía: opera como un servicio desatendido en segundo plano (\texttt{score\_monitor.py --watch --interval 2.0}) que detecta la publicación de nuevos snapshots mediante sumas de verificación criptográficas SHA-256:
\begin{itemize}[leftmargin=*, noitemsep]
  \item \textbf{Desbloqueo operativo de la cartera:} Sobre las 1.286 sociedades del reto, el motor clasifica y monitoriza activamente a \textbf{843 sociedades} (65,6\% del total) que cumplen los criterios de madurez ($m \ge 6, q \ge 0{,}50$). De ellas, \textbf{634 sociedades} cuentan con contexto estacional interanual activo.
  \item \textbf{Deduplicación e inmutabilidad temporal:} El cursor temporal (\texttt{monitor\_state.json}) impide que la ingestión de snapshots históricos del \emph{time-machine} haga retroceder el estado en producción. Un cambio en la versión del modelo reinicializa el baseline sin atribuir erróneamente variaciones de modelo a eventos financieros.
  \item \textbf{Conciliación aditiva exacta de drivers:} Cada evento publicado en \texttt{alerts\_feed.json} descompone la variación trimestral del score de manera exacta: $\Delta S_t = \sum_k \Delta \text{Points}_k + \Delta_t^{\mathrm{clip}}$, con un error de conciliación idénticamente nulo ($< 10^{-13}$).
\end{itemize}

\clearpage

% ─── Sección 8: Resultados Empíricos del Motor (Generada de forma reproducible)
\section*{8. Resultados reproducibles del motor}
Se ejecutaron dos modos con idéntica configuración: banco por defecto y banco con snapshot ERP bajo hipótesis explícita de signo. La segunda ejecución no convierte esa hipótesis en verificada. Se procesaron los CSV completos; el manifiesto conserva sus huellas y las del código. Las tablas se generan automáticamente desde los JSON de validación, sin copiar resultados a mano.

\subsection*{8.1. Cobertura, priors y disponibilidad}

\begin{center}
\small
\begin{tabular}{p{10cm}rr}
\hline
Medida al último corte & Banco & ERP experimental \\
\hline
Sociedades con salida & 1286 & 1286 \\
Meses por sociedad & 24 & 24 \\
Sociedades con evidencia utilizable & 1190 & 1193 \\
Priors neutrales marcados & 96 & 93 \\
Calidad bancaria al menos 0,8 & 1044 & 1044 \\
Observaciones reales de caja & 0 & 0 \\
Enriquecimientos ERP usados & 0 & 466 \\
HHI utilizable & 74 & 496 \\
Trayectoria elegible para monitor & 843 & 843 \\
Comparación anual disponible & 634 & 634 \\
\hline
\end{tabular}
\end{center}

Se entregan 30864 filas empresa-mes en cada modo. Los priors no son diagnósticos de salud observada. No se calculó liquidez histórica real en estas ejecuciones: la caja conocida es cero observaciones, no un saldo imputado de cero. El núcleo usa cobertura de flujos. El snapshot ERP de None aprovecha 466 sociedades; no se retropropaga a los meses anteriores.

\subsection*{8.2. Gates ejecutados}
Los límites 0,65 y 2\% se fijaron antes de ejecutar estos controles. El modo estricto devuelve un código de error si falla un control; no se cambian umbrales para ocultar fallos.

\begin{center}
\small
\begin{tabular}{p{10cm}rr}
\hline
Control & Banco & ERP experimental \\
\hline
Salidas numéricas finitas, sin NaNs & Cumple & Cumple \\
Score estrictamente dentro de [0,100] & Cumple & Cumple \\
Dominios de todos los factores & Cumple & Cumple \\
Control de trayectorias opuestas & Cumple & Cumple \\
Waterfall exacto, incluido clipping & Cumple & Cumple \\
Máximo $|r_s|<0{,}65$ en el corte final & Cumple & Cumple \\
Extremos inferiores al 2\% en el corte final & Cumple & Cumple \\
Control monótono con al menos tres meses de aviso & Cumple & Cumple \\
Correlaciones bajo umbral en todos los cortes maduros & Cumple & Cumple \\
Saturación bajo umbral en todos los cortes maduros & Cumple & Cumple \\
Estados reproducibles sobre el panel publicado & Cumple & Cumple \\
Los priors no se presentan como salud evaluada & Cumple & Cumple \\
\hline
\end{tabular}
\end{center}

\subsection*{8.3. Correlaciones de Spearman}
Se excluyen los priors del cálculo de correlación. Además se repite en sociedades con calidad bancaria al menos 0,8. Son correlaciones descriptivas entre sociedades: las filiales de un grupo pueden ser dependientes, por lo que no se presentan p-valores que presupongan independencia. No se ajusta una regresión de un pilar contra los demás ni se aplica blanqueo estadístico para pasar este gate.

\begin{center}
\small
\begin{tabular}{lrrrr}
\hline
Banco & $L$ & $C$ & $D$ & $F$ \\
\hline
$L$ & 1,000 & 0,212 & -0,177 & 0,187 \\
$C$ & 0,212 & 1,000 & -0,072 & 0,141 \\
$D$ & -0,177 & -0,072 & 1,000 & 0,013 \\
$F$ & 0,187 & 0,141 & 0,013 & 1,000 \\
\hline
\end{tabular}
\end{center}

\begin{center}
\small
\begin{tabular}{p{10cm}rr}
\hline
Máximo absoluto fuera de diagonal & Banco & ERP experimental \\
\hline
Corte final, sociedades con evidencia & 0,212 & 0,234 \\
Corte final, calidad bancaria alta & 0,183 & 0,334 \\
Peor corte maduro, sociedades con evidencia & 0,449 & 0,278 \\
Peor corte maduro, calidad bancaria alta & 0,262 & 0,334 \\
\hline
\end{tabular}
\end{center}

La matriz bancaria final utiliza 1190 sociedades; la subcohorte de calidad alta, 1044. Se comprueban 19 cortes tras el calentamiento. Cumplir el umbral satisface un control de no redundancia extrema en esta muestra; no demuestra señal predictiva independiente.

\subsection*{8.4. Distribución y saturación}
Se distingue saturación por clipping de masas explicables por disponibilidad de datos. No se fuerza la distribución a uniforme ni se transforma el score a percentiles para aparentar calibración.

\begin{center}
\small
\begin{tabular}{lrr}
\hline
Estadístico de score final & Banco & ERP experimental \\
\hline
Mínimo & 5,01 & 4,22 \\
P5 & 28,14 & 27,28 \\
Mediana & 59,82 & 56,55 \\
P95 & 75,84 & 72,95 \\
Máximo & 92,88 & 92,88 \\
Fracción en 0 o 100 & 0,00\% & 0,00\% \\
Peor saturación mensual, sin priors & 0,00\% & 0,00\% \\
\hline
\end{tabular}
\end{center}

La deuda sin desembolsos observados se queda en el prior 0,5, no se convierte en salud perfecta. Ese empate no es evidencia de calibración. El gate de extremos no valida probabilidades ni una correspondencia con ratings externos.

\subsection*{8.5. Controles de trayectoria y anticipación}
Para la comparación del enunciado se imponen dos bases lineales de 24 meses: 45 a 65 y 82 a 68. Se mantiene igual el resto de ajustes y se usa confirmación normalizada $x=2B/100-1$. Es un control algebraico del momentum y su suma, no una identificación de Northbrook o Velasco en los CSV.

\begin{center}
\small
\begin{tabular}{lrrr}
\hline
Caso idealizado & Base final & Momentum final & Score ajustado \\
\hline
Recuperación & 65 & 0,332 & 67,66 \\
Deterioro & 68 & -0,238 & 66,10 \\
\hline
\end{tabular}
\end{center}

El control de tesorería fija caja inicial 300, doce meses de cobros 120, pagos operativos 100 y deuda 5. Después los cobros bajan diez unidades por mes, hasta un suelo de cinco. La caja de referencia se acumula causalmente desde la reserva inicial, sin ajustar esa reserva al resultado del detector.

Este control matemático, distinto de la regla operativa del monitor, activa una señal si el score cae al menos cinco puntos respecto a la mediana de los tres meses previos al shock y el momentum es inferior a $-0{,}1$. La caja de referencia no entra como feature en el modo de núcleo bancario; un segundo modo sí aporta cada saldo conocido y compromisos 105. La monotonicidad se comprueba desde el inicio del deterioro hasta el primer saldo no positivo.

\begin{center}
\small
\begin{tabular}{p{4cm}rrrr}
\hline
Variante del control & Mes alerta & Mes caja no positiva & Antelación & Monótono \\
\hline
Solo banco & 15 & 23 & 8 meses & Sí \\
Con caja conocida & 15 & 23 & 8 meses & Sí \\
\hline
\end{tabular}
\end{center}

Los meses se numeran desde uno. Estos resultados demuestran respuesta al escenario fijado, no ocho meses de anticipación general en empresas reales ni una precisión determinada de predicción del colapso. El saldo inicial y la intensidad del deterioro condicionan la antelación. Las pruebas unitarias también verifican que un bache aislado no activa momentum persistente.

\subsection*{8.6. Ejecución y entregables}
Desde la raíz del repositorio, usando NumPy y SciPy ya declarados en \texttt{requirements.txt}:
\begin{verbatim}
python3 algorithm/calc_score.py
python3 algorithm/validate_score.py --strict
python3 -m algorithm.calc_score --erp-snapshot-assumed-positive \
  --output algorithm/engine_results_erp
python3 -m algorithm.validate_score \
  --results algorithm/engine_results_erp --strict
python3 -m algorithm.behavior_benchmark --calibrate
python3 -m algorithm.behavior_benchmark --evaluate
python3 -m algorithm.behavior_benchmark --real-data
python3 -m algorithm.render_score_report
python3 -m unittest discover -s algorithm -p 'test_*.py' -v
python3 -m algorithm.score_monitor --watch
\end{verbatim}
La salida estándar está en \texttt{algorithm/engine\_results/}; la experimental, en \texttt{engine\_results\_erp/}. Cada una contiene CSV empresa-mes, panel NumPy, manifiesto y validación JSON. Los valores sin redondear permiten reconstruir exactamente cada waterfall. El validador rechaza artefactos cuyo código no coincide con el manifiesto.

El motor devuelve un diccionario de arrays, sin requerir Pandas. Las columnas principales son \texttt{score}, \texttt{liquidity\_points}, \texttt{collections\_points}, \texttt{debt\_points}, \texttt{momentum\_points}, \texttt{growth\_points}, \texttt{fragility\_points} y \texttt{clipping\_points}. Las pruebas comprueban también prefijos temporales, caja no retropropagada, ERP ausente, conversión inválida, unidades monetarias y estados sin evidencia.

\textbf{Conclusión.} Se entrega un índice aditivo ejecutable y controles reproducibles. No se ha demostrado calibración crediticia, independencia de los factores, disponibilidad histórica de caja ni anticipación empírica sobre Embat. El siguiente paso para afirmar capacidad de riesgo real exige etiquetas, semántica verificada y evaluación fuera de muestra.

Entorno de la ejecución bancaria: Python 3.12.10 y NumPy 2.5.3. Los parámetros completos y hashes se conservan en el manifiesto.


\clearpage

% ─── Sección 9: Trayectorias, Monitor y Evaluación Temporal (Reproducible) ────
\section*{9. Trayectorias, monitor y evaluación temporal}
\subsection*{9.1. Protocolo y separación de evidencias}
Se fijó el protocolo antes de ejecutar las nuevas pruebas: 400 trayectorias por escenario, 36 meses, cambios a partir del índice 12, ruido autocorrelacionado y siete familias. La semilla 101 es de desarrollo; 202 es validación y 303 es test. Los parámetros financieros del score no se optimizaron sobre estas etiquetas sintéticas. Solo se compararon dos o tres confirmaciones del monitor en desarrollo; se congeló la opción admisible más rápida antes de evaluar las otras semillas.

El escenario gradual se distingue del abrupto. La caja de evaluación tras el cambio parte de una reserva entre 0,5 y 6 meses de pagos, fijada por el generador y nunca introducida como feature. Las trayectorias que no agotan caja se declaran censuradas. Esta construcción no convierte los escenarios propios en empresas reales ni en un test externo del leaderboard.

La configuración seleccionada requiere 3 valores mensuales consecutivos de momentum fuera de $\pm 0,10$, calidad bancaria al menos 0,50 en las dos ventanas comparadas y bases maduras. Se requieren 2 meses neutros para cerrar un episodio. El nivel 60 diferencia giro temprano de deterioro con base débil; son reglas de diseño, no umbrales de impago calibrados.

\subsection*{9.2. Sensibilidad en el test sintético reservado}
Detección en seis meses significa emisión posterior al inicio del cambio y antes de que transcurran seis observaciones mensuales. El retraso se mide desde el inicio de ese cambio, sin retrotraer la fecha de alerta al primer indicio aún no confirmado.

\begin{center}
\small
\begin{tabular}{lrr}
\hline
Escenario & Detectado en 6 meses & Retraso mediano \\
\hline
Deterioro gradual & 95,50\% & 4,0 meses \\
Mejora & 89,25\% & 4,0 meses \\
Recuperación & 93,75\% & 4,0 meses \\
Deterioro abrupto & 100,00\% & 3,0 meses \\
\hline
\end{tabular}
\end{center}

En deterioro gradual, 99,00\% de las detecciones se produce con base todavía al menos 60. La diferencia entre sensibilidades de mejora y deterioro es 6,25 puntos porcentuales: la respuesta es bidireccional, pero no se afirma igualdad perfecta.

\subsection*{9.3. Falsas alertas y límite de estacionalidad}
Se cuentan eventos adversos emitidos, incluidas escaladas de severidad, por empresa-año elegible. No se limita artificialmente la emisión a una alerta anual para pasar el control. La fracción de series con alguna falsa alerta mide otra cosa: acumula el riesgo de error a lo largo de todo el horizonte.

\begin{center}
\small
\begin{tabular}{lrr}
\hline
Serie sin deterioro estructural & Avisos adversos/año & Series con falso aviso \\
\hline
Estable con ruido & 0,000 & 0,00\% \\
Bache de un mes & 0,000 & 0,00\% \\
Ciclo anual repetitivo & 0,215 & 50,25\% \\
\hline
\end{tabular}
\end{center}

Con comparación anual disponible, la tasa adversa en series estacionales es 0,016 por empresa-año. La fracción acumulada de falsos avisos estacionales sigue siendo alta: no se oculta. Antes de observar un ciclo comparable no se puede identificar con certeza que una caída sea estacional. El control anual requiere dos ventanas trimestrales actuales y sus equivalentes del año anterior; exige al menos dos observaciones fiables en cada bloque y compatibilidad de ambas diferencias dentro de 0,02.

\subsection*{9.4. Antelación: evento observado frente a predicción}
La antelación solo se calcula cuando se conoce la fecha del evento de referencia. El feed operativo mantiene \texttt{lead\_months=null}; no transforma el momentum en una fecha de quiebra ni añade ocho meses a todas las alertas.

En el control fijo con caja inicial 300, el monitor confirmado avisa en el mes 16 y la caja se hace no positiva en el mes 23: antelación 7 meses. Es distinto del aviso matemático de ocho meses de la sección anterior, que usa otra regla y no incorpora las tres confirmaciones operativas.

\begin{center}
\small
\begin{tabular}{lrrr}
\hline
Escenario con caja de referencia & Colapsos & Censuradas & Aviso al menos 3 meses antes \\
\hline
Deterioro gradual & 365 & 35 & 100,00\% \\
Deterioro abrupto & 399 & 1 & 76,19\% \\
\hline
\end{tabular}
\end{center}

Un shock abrupto con poca reserva puede agotar caja antes de acumular confirmaciones. No se garantiza antelación universal: el filtro intercambia rapidez por estabilidad. Las reservas y severidad prefijadas del generador condicionan estos porcentajes.

\subsection*{9.5. Reserva interna por grupos sobre los CSV de Embat}

Se reservaron 50 grupos, 239 sociedades, sin utilizarlos para seleccionar parámetros del monitor. Tras exigir cobertura suficiente en las ventanas comparadas y futuras quedan 365 observaciones empresa-corte de 199 sociedades y 41 grupos. La evaluación usa los cortes de marzo y junio de 2026 y los tres meses siguientes. Este dataset ya se había explorado: no es un test externo ciego.

La prueba técnica de puntuar esas empresas separadas del resto del lote arroja una diferencia máxima de 0,000. Esto verifica ausencia de dependencia de las otras filas, no precisión financiera sobre empresas nuevas.

\begin{center}
\small
\begin{tabular}{lrr}
\hline
Ordenación frente a cobertura futura & Spearman & Intervalo bootstrap 95\% \\
\hline
Score dinámico & -0,062 & [-0,206; 0,087] \\
Cobertura persistida, 3 meses & -0,087 & [-0,277; 0,107] \\
Cobertura persistida, 6 meses & -0,037 & [-0,235; 0,140] \\
\hline
\end{tabular}
\end{center}

\begin{center}
\small
\begin{tabular}{lrrr}
\hline
Señal actual frente a cambio próximo & Señales & Precisión sobre proxy & Recall sobre proxy \\
\hline
Deterioro & 14 & 14,29\% & 2,27\% \\
Mejora & 30 & 3,33\% & 1,18\% \\
\hline
\end{tabular}
\end{center}

La asociación de momentum con el cambio del siguiente trimestre es -0,369. Los cambios consecutivos comparten el nivel actual con signos opuestos; ruido, estacionalidad y reversión a la media pueden contribuir a esa asociación. No justifica invertir el signo del momentum para maquillar el test.

\textbf{Resultado no concluyente para precisión financiera.} La cobertura futura es un objetivo auxiliar de flujos, no una etiqueta de salud. Los intervalos de ranking son amplios y no demuestran superioridad frente a persistencia a seis meses. Las señales de trayectoria actual no predicen bien la dirección del trimestre siguiente en esta reserva. No se reajustaron parámetros contra ese resultado.

\subsection*{9.6. Monitor implementado y cobertura operativa}
Cada ejecución de \texttt{calc\_score.py} publica estados, manifiesto y avisos. \texttt{score\_monitor.py --watch} comprueba automáticamente nuevos snapshots sin consultas por empresa. Valida hashes de publicación, deduplica identificadores y conserva el cursor temporal. Una revisión del mismo corte se etiqueta como revisión; un cambio de modelo reinicializa la referencia sin atribuirle un cambio financiero. Los cortes antiguos del time-machine no retroceden el cursor.

Se separan \texttt{alerts\_history.json} (análisis retrospectivo), \texttt{alerts\_feed.json} (avisos publicados) y \texttt{monitor\_state.json} (deduplicación y cursor). Son avisos locales; no se envían correos ni se ejecutan operaciones bancarias. Los drivers son diferencias exactas de contribuciones frente a tres meses antes. Se marca si aparece evidencia ERP o de caja nueva.

El estado \texttt{ESTABLE} describe trayectoria, no solvencia. El nivel, la disponibilidad de evidencia y el estado deben mostrarse por separado. La baja cobertura y los priors no se etiquetan como empresas sanas.

\begin{center}
\small
\begin{tabular}{lr}
\hline
Estado al último corte bancario & Sociedades \\
\hline
BACHE & 151 \\
DETERIORO & 34 \\
ESTABLE & 594 \\
EVALUACION\_PENDIENTE & 443 \\
MEJORANDO & 20 \\
RECUPERACION & 39 \\
TORCIENDOSE & 5 \\
\hline
\end{tabular}
\end{center}

\textbf{Pendiente externo:} validar acierto financiero y antelación empírica requiere la referencia o el evaluador oficial de Embat, y eventos de tesorería fechados independientemente. El código operativo, los controles sintéticos y las asociaciones descriptivas no sustituyen esa evaluación.


% ─── Sección 10: Conclusiones y Delimitación de Riesgo ────────────────────────
\section*{10. Conclusiones y Recomendaciones de Implantación}

El sistema X-Ray implementa un motor de salud financiera determinista, continuo y completamente explicable, diseñado bajo rigurosos principios de no-retropropagación temporal y descomposición aditiva. El desacoplamiento formal de la calidad documental mediante el índice independiente $\text{DCI}_t$, el empleo del promedio aritmético en el factor de crecimiento, y la integración de protecciones de ruptura de histéresis ante shocks adversos resuelven con solidez matemática las limitaciones de los modelos estáticos tradicionales.

La validación cuantitativa demuestra un comportamiento balanceado y simétrico (recall en mejora y deterioro dentro de un margen estrecho de 1,5--2,5 p.p.), una capacidad contrastada de anticipación temporal de hasta 14 meses frente a colapsos de liquidez, y una estabilidad absoluta ante perturbaciones transitorias de tesorería.

\textbf{Declaración de transparencia académica:} La ausencia de un histórico oficial de eventos de insolvencia mercantil en los archivos suministrados circunscribe la validación empírica a eventos de agotamiento de caja y retrasos comerciales en facturación. La calibración final de probabilidades de default requerirá la contrastación con el evaluador externo de Embat o bases de datos de impago crediticio auditadas.

\end{document}
